El análisis de la cartera de crédito muestra que, de un total de 1,000 clientes, 127 presentan incumplimiento en sus obligaciones de pago, lo que representa una tasa de mora del
12.7%
. En contraste, el
87.3%
de los clientes mantiene un comportamiento de pago al día, reflejando que la mayor parte de la cartera no presenta atrasos.
Aunque la mayoría de los clientes cumple con sus obligaciones financieras, una tasa de mora del
12.7%
constituye un indicador relevante que debe evaluarse en función de las políticas internas de riesgo y de los objetivos establecidos por la institución financiera. Su interpretación dependerá del nivel máximo de mora considerado aceptable por el banco.
Desde una perspectiva de gestión del riesgo, resulta recomendable profundizar el análisis para identificar las características asociadas a los clientes en situación de mora, considerando variables como el nivel de ingresos, el monto del préstamo, la antigüedad laboral, el sector económico y la edad. Este tipo de análisis permitirá diseñar estrategias de prevención, fortalecer los procesos de otorgamiento de crédito y reducir el riesgo de incumplimiento en el futuro.